%% ma: L10-B13-0010
%% lop: 10
%% chuong: 5
%% bai: 13
%% dang: 10.13.4
%% loai: TL
%% muc: TH
%% dap_an: "-"
%% nguon: "chuyen-de-cac-so-dac-trung-cua-mau-so-lieu-khong-ghep-nhom-toan-10-chuong-trinh-moi.pdf (GV. Phan Nhật Linh) (P02), Bài 02 Dạng 1 Bài tập tự luận 10"
%% kien_thuc: "Lớp 10 Bài 13 (trung vị)"
%% trang_thai: chinh-thuc
%% ghi_chu: "Tự duyệt theo yêu cầu giáo viên 22/10/2026."
\begin{ex}
Số áo bán được của một cửa hàng được cho trong bảng sau:
\begin{center}
\begin{tabular}{|c|c|c|c|c|c|c|c|}\hline
Cỡ áo & $36$ & $37$ & $38$ & $39$ & $40$ & $41$ & $42$\\\hline
Tần số & $33$ & $75$ & $126$ & $80$ & $106$ & $40$ & $5$\\\hline
\end{tabular}
\end{center}
Tìm trung vị của mẫu số liệu.
\loigiai{$n=465$ (lẻ) nên trung vị là giá trị thứ $233$.

Tần số tích lũy: $33$; $108$; $234$ nên giá trị thứ $233$ là $38$. Vậy $M_e=38$.}
\end{ex}
